\documentclass[12pt]{article}
\usepackage[utf8]{inputenc}
\usepackage[T1]{fontenc}
\usepackage{hyperref}
\title{Eastern Yiddish Language Resources for K-12 Curriculum Development}
\author{Compiled by Claude Sonnet 5.5 (claude-sonnet-5-5) for Andr\'as Kornai}
\date{}
\begin{document}
\maketitle

\section{Introduction}
Eastern Yiddish (ISO 639-3 ydd) is spoken by about one million people, and is written in the Hebrew alphabet. The search found several textbooks and courses, mostly for adults and universities, plus some children's materials from Chabad-oriented and historical sources. This is a short survey from one general search. It is not exhaustive.

\section{Primary Educational Resources}
\begin{itemize}
\item Sheva Zucker, \emph{Yiddish: An Introduction to the Language, Literature and Culture}, Vol.\ 1, a widely used adult textbook with an answer key and recordings. \url{https://shop.yiddishbookcenter.org/products/yiddish-an-introduction-to-the-language-literature-and-culture-vol-1-book-by-zucker-sheva}
\item Mordkhe Schaechter, \emph{Yiddish II: An Intermediate and Advanced Textbook}, League for Yiddish. \url{https://leagueforyidev.leagueforyiddish.org/en/bikher/yidish2}
\item Chabad curriculum database, lower elementary Yiddish; it includes a workbook for grades 3--8 that teaches Yiddish from English. \url{https://mymef.org/chabad-curriculum-database/lower-elementary/yiddish/}
\item Chabad curriculum database, upper elementary Yiddish. \url{https://mymef.org/chabad-curriculum-database/upper-elementary/yiddish/}
\end{itemize}
Historical school readers exist, such as Steinbaum's \emph{Leyenbikher far der idisher shul} (New York, 1924), but I did not find a current source for them.

\section{Linguistic and Reference Materials}
No dictionary or grammar was specifically searched for or verified in this pass.

\section{Religious and Community Resources}
\begin{itemize}
\item YIVO Institute, intensive beginner Yiddish course (Fall 2022). \url{https://yivo.org/FALL2022-Intensive-Beginner1and2}
\item League for Yiddish (publisher of the Schaechter textbook, above).
\end{itemize}

\section{Digital and Audio Resources}
\begin{itemize}
\item Free beginner Yiddish textbook audio lessons online. \url{https://freelanguage.org/node/88005}
\end{itemize}

\section{Curriculum Development Context}
Yiddish courses exist at universities, including the University of Antwerp (level 1) and Berkeley, and in the Chabad curriculum for elementary grades. Search results did not show a full K-12 curriculum conducted in Yiddish, so building one would need adapting adult and higher-education materials and the elementary workbooks listed above.
\begin{itemize}
\item University of Antwerp, Yiddish level 1. \url{https://www.uantwerpen.be/en/rg/ijs/courses/yiddish/level-1/}
\end{itemize}
\end{document}